% GAL1897_W05 — English translation of physical PDF pages 040–041.
% Primary authority: the frozen 1897 French diplomatic transcription from
% JP2 leaves L0039–L0040, navigated against the 96-page 1897 PDF.
% Definite source-level mathematical errors are translated without silent repair
% and identified in adjacent critical notes from the frozen critical baseline.

\providecommand{\GalSourcePageStart}[4]{}%
\providecommand{\GalSourcePageBoundary}[4]{\space}%
\providecommand{\GalSourceError}[2]{#1}%
\providecommand{\GalHistoricalFootnote}[1]{\footnote{#1}}%
\providecommand{\GalArticleEndRule}{\par\smallskip\begin{center}\rule{0.18\linewidth}{0.4pt}\end{center}}%
\providecommand{\GalOrdinal}[1]{#1\textsuperscript{o}}%

\GalSourcePageStart{040}{L0039}{11}{11}

% ENSEG:W05:PDF040:0001
\begin{center}
{\Large\bfseries ANALYSIS\par}
\vspace{0.2em}
{\scriptsize\bfseries OF A\par}
\vspace{0.2em}
{\normalsize\bfseries MEMOIR ON THE ALGEBRAIC SOLUTION OF EQUATIONS\textsuperscript{(*)}.\par}
\vspace{0.75em}
\rule{0.12\linewidth}{0.4pt}
\end{center}

% ENSEG:W05:PDF040:0002
\GalHistoricalFootnote{M.~Férussac's \emph{Bulletin des Sciences mathématiques}, vol.~XIII, p.~271
(1830, April issue). \textsc{(J. Liouville.)}}

% ENSEG:W05:PDF040:0003
Equations are called nonprimitive when, being, for example, of
degree $mn$, they decompose into $m$ factors of degree $n$ by means
of a single equation of degree $m$. These are M.~Gauss's equations.
Primitive equations are those that admit no such simplification. With
regard to primitive equations, I have obtained the following results:

\medskip
% ENSEG:W05:PDF040:0004
\noindent\GalOrdinal{1} For an equation of prime degree to be solvable by
radicals, it is necessary and sufficient that, once any two of its roots
are known, the others can be deduced rationally from them.

\smallskip
% ENSEG:W05:PDF040:0005
\noindent\GalOrdinal{2} For a primitive equation of degree $m$ to be solvable
by radicals, it is necessary that $m=p^{\nu}$, where $p$ is a prime number.

\smallskip
% ENSEG:W05:PDF040:0006
\noindent\GalOrdinal{3} \GalSourceError{Apart from the cases mentioned below, for a
primitive equation of degree $p^{\nu}$ to be solvable by radicals, it is
necessary that, once any two of its roots are known, the others can be
deduced rationally from them.}{W05-SE001: the printed exception condition is too restrictive; preserve the wording and consult the separate degree-8 affine-semilinear counterexample certificate.}

% ENSEG:W05:PDF040:0007
\GalCriticalNote{W05-SE001}{A degree-$16$ counterexample to the exception list}{The
finite list of exceptions is not a complete modern classification. On
$V=\mathbf F_2^4$, an irreducible solvable complement
$(C_3\times C_3)\rtimes C_2$ of order $18$ gives a primitive affine group of
order $288$. It is not one-dimensional semilinear, since
$18\nmid|\Gamma L(1,16)|=60$. In general, a finite solvable primitive group is
affine, with a regular elementary-abelian socle and a faithful irreducible point
stabilizer; one-dimensional semilinearity requires an additional hypothesis.
The historical sentence is left unchanged.}

\medskip
% ENSEG:W05:PDF040:0008
The following very special cases escape the preceding rule:

\smallskip
% ENSEG:W05:PDF040:0009
\noindent\GalOrdinal{1} The case $m=p^{\nu}=9$, $=25$;

\smallskip
% ENSEG:W05:PDF040:0010
\noindent\GalOrdinal{2} The case $m=p^{\nu}=4$, and generally the case in which,
with $a^{\alpha}$ a divisor of $\dfrac{p^{\nu}-1}{p-1}$, $a$ is prime and one has
% ENSEG:W05:PDF040:0011
\[
\frac{p^{\nu}-1}{a^{\alpha}(p-1)}\,\nu
  =p\quad(\mathrm{mod.}\ a^{\alpha}).
\]
% ENSEG:W05:PDF040:0012
These cases nevertheless depart very little from the general rule.

% ENSEG:W05:PDF040:0013
When $m=9$, $=25$, the equation must be of the type of those that
determine the trisection and quintisection of elliptic functions.

\GalSourcePageBoundary{041}{L0040}{12}{12}

% ENSEG:W05:PDF041:0014
In the second case it will always be necessary that, once two of the
roots are known, the others can be deduced from them, at least by means
of as many radicals of degree $p$ as there are divisors $a^{\alpha}$ of
$\dfrac{p^{\nu}-1}{p-1}$ that satisfy
% ENSEG:W05:PDF041:0015
\[
\frac{p^{\nu}-1}{a^{\alpha}(p-1)}\,\nu
  =p\quad(\mathrm{mod.}\ a^{\alpha}),\qquad a\text{ prime}.
\]

% ENSEG:W05:PDF041:0016
All these propositions were deduced from the theory of permutations.

% ENSEG:W05:PDF041:0017
Here are further results that follow from my theory.

\medskip
% ENSEG:W05:PDF041:0018
\noindent\GalOrdinal{1} Let $k$ be the modulus of an elliptic function and $p$ a
given prime number $>3$. For the equation of degree $p+1$ that gives the
various moduli of the functions transformed with respect to the number
$p$ to be solvable by radicals, \emph{one of two things is necessary}:
either one of the roots must be known rationally, or all the roots must
be rational functions of one another. It is of course understood that
only particular values of the modulus $k$ are at issue here. It is clear
that this does not hold in general. This rule does not hold for $p=5$.

\medskip
% ENSEG:W05:PDF041:0019
\noindent\GalOrdinal{2} It is remarkable that the general modular equation of
the sixth degree corresponding to the number $5$ can be reduced to one
of the fifth degree of which it is the reduced equation.
\GalSourceError{By contrast, in higher degrees the modular equations cannot
be reduced in degree}{W05-SE002: the printed impossibility is false for p=7 and p=11, as the 1897 Liouville footnote itself records; preserve the wording and consult the separate group-theoretic certificate}\textsuperscript{(*)}.

% ENSEG:W05:PDF041:0020
\GalCriticalNote{W05-SE002}{Higher-degree reductions do exist}{The main text
asserts a universal impossibility in higher degrees, but Liouville's following
note records reductions for $p=7$ and $p=11$ as well as $p=5$. Exact subgroup
calculations confirm the $p=7$ and $p=11$ exceptions. For $p=11$, the corrected
six-pair matching has stabilizer order $60$ and index $11$ in $PSL(2,11)$.
Liouville's printed 1846 footnote is the earliest explicit correction found;
the erroneous sentence and its historical correction are kept separate here.}

% ENSEG:W05:PDF041:0021
\GalHistoricalFootnote{This assertion is not entirely accurate, as Galois himself warns
in his Letter to M.~Auguste Chevalier, which appears below. Speaking generally
of the article reproduced here, he says: ``The condition I gave in the
\emph{Bulletin de Férussac} for solvability by radicals is too restrictive;
there are few exceptions, but there are some.'' As for modular equations in
particular, he declares that the reduction from degree $p+1$ to degree $p$ is
possible not only for $p=5$, but also for $p=7$ and $p=11$; he nevertheless
maintains its impossibility for $p>11$. \textsc{(J. Liouville.)}\GalArticleEndRule}
